% !TEX program = xelatex
% ReSETP full-paper evidence-visible draft, 2026-08-19.
\documentclass{setp-new}

\usepackage{fontspec}
\setmainfont{Times New Roman}
\setsansfont{Arial}
\IfFontExistsTF{Songti SC}{\newcommand{\setpSongFont}{Songti SC}}{\newcommand{\setpSongFont}{SimSun}}
\IfFontExistsTF{Songti SC}{\newcommand{\setpFangSongFont}{Songti SC}}{\newcommand{\setpFangSongFont}{FangSong}}
\IfFontExistsTF{SimHei}{\newcommand{\setpHeiFont}{SimHei}}{%
  \IfFontExistsTF{Noto Sans CJK SC}{\newcommand{\setpHeiFont}{Noto Sans CJK SC}}{\newcommand{\setpHeiFont}{\setpSongFont}}}
\setCJKmainfont[BoldFont=\setpHeiFont]{\setpSongFont}
\setCJKsansfont{\setpHeiFont}
\setCJKmonofont{\setpSongFont}
\setCJKfamilyfont{zhsong}[BoldFont=\setpHeiFont]{\setpSongFont}
\setCJKfamilyfont{zhhei}{\setpHeiFont}
\setCJKfamilyfont{zhfs}{\setpFangSongFont}
\providecommand{\songti}{}
\providecommand{\heiti}{}
\providecommand{\fangsong}{}
\renewcommand{\songti}{\CJKfamily{zhsong}}
\renewcommand{\heiti}{\CJKfamily{zhhei}}
\renewcommand{\fangsong}{\CJKfamily{zhfs}}
\XeTeXlinebreaklocale "zh"
\XeTeXlinebreakskip=0pt plus 1pt

\usepackage{mathtools}
\usepackage{longtable}
\usepackage{graphicx}
\usepackage{pgfplots}
\usepackage{pgfplotstable}
\usetikzlibrary{arrows.meta,positioning,fit,calc,patterns}
\pgfplotsset{compat=1.18}
\setlength{\parindent}{2em}
\setlength{\headheight}{32pt}
\setlength{\emergencystretch}{3em}
\setlength{\LTpre}{4pt}
\setlength{\LTpost}{4pt}

% 证据可见占位。正文中所有尚无合格数据的数字和结论均使用此宏。
\DeclareRobustCommand{\TBD}[1]{%
  \begingroup\setlength{\fboxsep}{1.2pt}%
  \colorbox{black!12}{\textcolor{black!65}{\textsf{\footnotesize[待填：#1]}}}%
  \endgroup}
\newcommand{\DraftNote}[1]{%
  \par\smallskip\noindent\fcolorbox{black!25}{black!5}{%
    \parbox{0.95\linewidth}{\small\textcolor{black!70}{#1}}}\par\smallskip}
\newcommand{\sourcelevel}{\textsuperscript{\scriptsize 数据身份见表注}}
\newcommand{\yes}{是}
\newcommand{\no}{否}

\jolname{系统工程理论与实践}{Systems Engineering --- Theory \& Practice}
\Volume{20XX}{1}{XX}{X}
\PageNum{1}
\PaperDOI{}
\Category{}{A}

\TitleMark{时变电网碳强度下的混合车队动态配送}
\Title{时变电网碳强度下的混合车队动态配送路径优化模型及算法}
\Author{\TBD{作者姓名}}{\TBD{作者单位、城市和邮政编码}}
\Abstract{%
在交通运输低碳转型背景下，电动配送车辆的运行阶段虽不产生尾气排放，但其充电侧排放取决于充电发生时段的电网碳强度。本文提出时变电网碳强度下的动态多车场混合车队多趟车辆路径问题，在统一框架中优化客户服务、车型指派、实体车多趟排班、充电时刻与动态发车决策，并用燃油车直接排放与电动车充电侧间接排放之和度量系统全口径排放。求解方面，本文在成熟混合遗传搜索内核上组织一条问题特化的串行改进链，使完整成本与完整可行性控制候选的接受、入群与生存。固定排班的作用空间计算显示，碳感知择时相对有空即充可移动充电量为174.732~kWh，充电侧排放减少59.082~kgCO$_2$e，但在0.075元/kg的碳价下完整目标增加22.094元；两种策略的配对盈亏平衡碳价为0.449元/kg。协同效率批中，两家单干成本之和为3922.079元，联合成本为3033.206元，节约888.873元；事后分摊后两家各节约444.436元。公开算例的稳定排位、增量式消融、完整车型扫描与动态需求结果仍分别标为\TBD{对应正式结果}。}
\Keywords{时变电网碳强度；混合车队；动态需求；多车场；多趟车辆路径；混合遗传搜索}

\ETitle{\TBD{英文题名尚未批准}}
\EAuthor{\TBD{英文作者姓名}}{\TBD{英文作者单位与地址}}
\EAbstract{This draft formulates the dynamic multi-depot mixed-fleet multi-trip vehicle routing problem with time-varying grid carbon intensity. The model jointly determines customer service, vehicle-type assignment, physical-vehicle multi-trip schedules, charging times, and dynamic dispatch decisions. A serial problem-specific improvement chain is organized on a mature hybrid genetic search kernel, while full cost and full feasibility govern population decisions. In a fixed-schedule opportunity calculation, carbon-aware charging moves 174.732~kWh and reduces charging emissions by 59.082~kgCO$_2$e relative to charging as soon as possible, but increases the complete monetary objective by 22.094~CNY at a carbon price of 0.075~CNY/kg; the paired break-even price is 0.449~CNY/kg. In the collaboration batch, separate operation costs 3922.079~CNY and joint operation costs 3033.206~CNY, yielding 888.873~CNY of savings; post-hoc allocation gives each carrier 444.436~CNY of savings. \TBD{benchmark ranking, accepted ablation, full fleet scan, and dynamic-demand findings}.}
\EKeywords{time-varying grid carbon intensity; mixed fleet; dynamic demand; multi-depot; multi-trip vehicle routing; hybrid genetic search}

\begin{document}
\maketitle

\DraftNote{本文是“结构完整、真假分明”的可编译草稿，不是投稿定稿。灰底\TBD{占位}表示缺少合格结果、来源或用户决定；已填数字的完整精度、来源字段和派生算式登记在随稿报告中。}

\section{引言}

实现碳达峰、碳中和要求推动运输工具装备低碳转型，建设绿色高效交通运输体系\cite{ref:state-council-2021}。城市群与区域配送中，多家配送企业可能分别拥有车场、有限车队与责任客户，并通过共享车辆、任务和补能资源降低系统总成本。协作形成的节约如何在各企业之间划分，关系到各方是否愿意持续参与。

电动配送车辆在运行阶段不产生尾气排放，但充电行为将排放转移到电力生产侧。因而，同一笔充电电量在不同电网时段对应的间接排放并不相同。若将电动车直接记为零排放，或用单一日均排放因子核算全部充电量，调度行为与实际充电侧排放之间的联系会被压平。Li等研究了有序充电的减排价值\cite{ref:li-2024}，Cheng等在车辆可用时段、站端功率和电池状态约束下研究了碳感知充电\cite{ref:cheng-2022}。这类研究回答“何时充电”，但配送车辆的可充窗口同时受路线、硬时间窗和多趟排班约束。

混合车队路径研究关注燃油车与电动车在载重、续航、能耗和固定成本上的差异。Goeke和Schneider建立了燃油车---电动车混合车队路径模型\cite{ref:goeke-2015}；李得成等研究了带时间窗的混合车辆路径问题\cite{ref:li-decheng-2021}；陈婉茹等分析了碳交易下多中心混合车队的路径与速度联合优化\cite{ref:chen-wanru-2023}。当充电排放随时刻变化时，可用电动车名额是否转化为实际派遣，以及车型结构和充电时刻各自贡献多少，需要使用同一全口径排放公式辨析。

动态车辆路径研究关注信息逐步揭示条件下的实时调度。Gendreau等采用并行禁忌搜索处理实时车辆路径与调度\cite{ref:gendreau-1999}，姜广田等比较了多车型动态车辆路径的调整策略\cite{ref:jiang-2024}。在多车场多趟场景中，动态重规划除更新待服务客户外，还应继承车辆位置、最早可用时刻、剩余载重、剩余电量以及已经发生的成本和排放。Zhen等对带释放时间的多车场多趟问题进行了系统建模\cite{ref:zhen-2020}。

基于上述背景，本文研究时变电网碳强度下的动态多车场混合车队多趟车辆路径问题（Dynamic Multi-depot Mixed-fleet Multi-trip Vehicle Routing Problem with Time-varying Grid Carbon Intensity，GCI-DMM-VRP）。当前稿将研究内容组织为四部分：
\begin{enumerate}[label=（\arabic*）]
  \item 建立同时考虑时变电网碳强度、混合车队、多车场、实体车多趟与动态需求的统一模型，使可用充电窗口由排班内生决定；
  \item 在同一代表实例和同一排放核算下，分别观察充电时刻、车型指派和动态发车三个决策层；
  \item 在成熟混合遗传搜索内核上组织问题特化组件，并用公开算例、增量式消融和收敛过程分别检验；
  \item 将协同效率与收益分配分开：先比较单干与联合运营，再在同一联合方案上事后分摊节约。
\end{enumerate}

\TBD{系统文献检索核验“统一考虑上述因素”的研究缺口及引用页码}。本文的实验结论强度由后文真实数据决定，不在此预写公开排位、动态效果或未完成机制的方向。

\section{问题描述与模型}

\subsection{问题描述}

考虑由多家配送企业协作运营的区域配送系统。系统包含车场集合、客户集合和公共充电站集合。每家企业拥有一个车场、数量有限的燃油车和电动车，以及一组责任客户；每辆实体车具有唯一所属车场，可在一个运营日内完成一趟或多趟配送任务。客户具有需求量、服务时间和硬时间窗，每个客户由一辆车完整服务一次。

协作允许一家企业的车辆服务另一家企业的责任客户，并在全系统范围内重新组织路线与车型；企业车队上限和车辆归属仍然成立。燃油车排放按燃油消耗计算，电动车排放按各时段充电量与该时段电网碳强度的乘积计算。运营过程中，订单事件逐步揭示；每次重规划继承车辆当前状态，已执行路径前缀不可撤回。

模型采用以下基本假设：客户需求不可拆分；车辆满足载重、硬时间窗与运营时域约束；固定成本按实际启用的实体车计取；电动车允许部分充电；车场充电不设置并发上限，公共站容量按已登记资源处理；动态重规划不重复结算已经发生的成本和排放。非线性充电的实验解释按“是否主动在边际收益递减段前停止充电”组织，而不是要求每次充电都触及曲线拐点。

\subsection{符号说明}

\begin{table}[htbp]
  \centering
  \caption{主要集合、参数、状态与决策变量}
  \label{tab:symbols}
  \small
  \begin{tabular}{p{0.12\linewidth}p{0.21\linewidth}p{0.48\linewidth}p{0.13\linewidth}}
    \toprule
    类别 & 符号 & 含义 & 单位/域\\
    \midrule
    集合 & $D,N,S,T,K$ & 车场、客户、公共充电站、电网时段与实体车辆集合 & ---\\
    集合 & $K_d,K^E,K^F$ & 车场$d$所属车辆、电动车与燃油车集合 & ---\\
    集合 & $\Omega_k$ & 实体车$k$的可行日排班模式集合；一个模式可含多趟 & ---\\
    路网参数 & $d_{ij},t_{ij}$ & 节点$i$至$j$的距离与行驶时间 & km，min\\
    客户参数 & $q_i,[e_i,l_i],s_i$ & 客户需求、时间窗与服务时长 & kg，min\\
    车辆参数 & $Q_k,B_k,b_k^{\min}$ & 载重上限、电池容量与最低允许电量 & kg，kWh\\
    补能参数 & $P_s,C_s$ & 充电功率与公共站容量 & kW，个\\
    时序参数 & $\gamma_t,p^e_{s,t}$ & 时段$t$的电网碳强度与电价 & kgCO$_2$e/kWh，元/kWh\\
    成本参数 & $c_k^{\mathrm{fix}},c_k^{\mathrm{km}},p^F,p^{\mathrm{car}}$ & 固定成本、里程成本、油价和单位碳价 & 相应单位\\
    路线变量 & $x_{ijr},a_{ir}$ & 趟$r$是否使用弧$(i,j)$、是否服务客户$i$ & $0$--$1$\\
    状态变量 & $L_{ir},T_{ir},b_{ir}$ & 离开节点$i$时的载重、服务时刻和电量 & kg，min，kWh\\
    充电变量 & $Y_q,h_q,\Delta_q,y_{qt}$ & 会话电量、开始时刻、持续时间和分时段电量 & kWh，min\\
    模式变量 & $A_{ikp},D_{kp},F_{kp}$ & 客户覆盖、模式里程与燃油消耗 & ---，km，L\\
    模式变量 & $z_{kp},u_k$ & 是否选择模式、是否启用实体车 & $0$--$1$\\
    动态状态 & $o_k^\tau,t_k^\tau,\ell_k^\tau,b_k^\tau$ & 时刻$\tau$的车辆位置、可用时刻、余量和电量 & ---\\
    \bottomrule
  \end{tabular}
  \tabnote{注：本表保留进入正文模型的主要符号；参数数值与证据身份见表\ref{tab:parameters}。}
\end{table}

\subsection{能耗、充电与时变碳排放核算}

在给定行驶速度下，电动车逐弧电耗与燃油车逐弧油耗写为距离和载重的函数：
\begin{align}
e^E_{ij,k}&=(\alpha_{0k}^E+\alpha_{1k}^E L_{ir})d_{ij}, && k\in K^E,\label{eq:ev-energy}\\
f^F_{ij,k}&=(\alpha_{0k}^F+\alpha_{1k}^F L_{ir})d_{ij}, && k\in K^F.\label{eq:cv-fuel}
\end{align}
其中，最终稿仍需逐式核对系数的单位、实现值和文献来源\cite{ref:goeke-2015}。

设充电会话$q$发生于站点$s(q)$，功率为$P_{s(q)}$，其与电网时段$t$的重叠时长为
\begin{equation}
g_{qt}=\left[\min\{h_q+\Delta_q,\bar T_t\}-\max\{h_q,\underline T_t\}\right]_+,
\label{eq:charge-overlap}
\end{equation}
则时段充电量为$y_{qt}=P_{s(q)}g_{qt}/60$，并满足$\sum_{t\in T}y_{qt}=Y_q$。部分充电和公共站容量的处理参照相关电动车路径研究\cite{ref:froger-2022,ref:keskin-2016}。

实体车排班模式$p$的燃油车直接排放和电动车充电侧间接排放分别为
\begin{align}
E_{kp}^{F}&=\lambda^F\sum_{r\in\mathcal R_{kp}}\sum_{(i,j)\in r}f^F_{ij,k},\label{eq:cv-emission}\\
E_{kp}^{E}&=\sum_{q\in\mathcal Q_{kp}}\sum_{t\in T}\gamma_ty_{qt}.\label{eq:ev-emission}
\end{align}
系统全口径排放为
\begin{equation}
E^{\mathrm{sys}}=\sum_{k\in K}\sum_{p\in\Omega_k}(E_{kp}^{F}+E_{kp}^{E})z_{kp}.
\label{eq:system-emission}
\end{equation}
式\eqref{eq:system-emission}是后文各决策层共同使用的排放尺：移动充电开始时刻改变$y_{qt}$的时段分布，改变车型指派改变直接与间接排放的构成，动态重规划则改变车辆可用时刻和充电窗口。

\subsection{目标函数与约束}

运营成本由实体车固定成本、里程成本、燃油成本和充电电费构成：
\begin{equation}
C^{\mathrm{op}}=\sum_{k\in K}c_k^{\mathrm{fix}}u_k+
\sum_{k\in K}\sum_{p\in\Omega_k}\left(c_k^{\mathrm{km}}D_{kp}+p^FF_{kp}+
\sum_{q\in\mathcal Q_{kp}}\sum_{t\in T}p^e_{s(q),t}y_{qt}\right)z_{kp}.
\label{eq:operation-cost}
\end{equation}
固定成本按启用实体车计取：
\begin{equation}
u_k=\sum_{p\in\Omega_k}z_{kp},\qquad \sum_{p\in\Omega_k}z_{kp}\leq 1.
\label{eq:vehicle-use}
\end{equation}
采用单位碳价把排放纳入统一货币目标：
\begin{equation}
\min Z=C^{\mathrm{op}}+p^{\mathrm{car}}E^{\mathrm{sys}}.
\label{eq:objective}
\end{equation}

每个客户由一个实体车排班模式恰好覆盖：
\begin{equation}
\sum_{k\in K}\sum_{p\in\Omega_k}A_{ikp}z_{kp}=1,\qquad i\in N.
\label{eq:coverage}
\end{equation}
模式内部满足流守恒、载重与硬时间窗约束：
\begin{align}
\sum_jx_{ijr}&=\sum_jx_{jir}=a_{ir}, && i\in N,\label{eq:flow}\\
0\leq L_{ir}&\leq Q_k,\qquad L_{jr}\leq L_{ir}-q_j+M(1-x_{ijr}),\label{eq:load}\\
e_i a_{ir}&\leq T_{ir}\leq l_i+M(1-a_{ir}),\nonumber\\
T_{jr}&\geq T_{ir}+s_i+t_{ij}-M(1-x_{ijr}).\label{eq:time-window}
\end{align}
电动车电量沿弧传播，并仅在允许地点补能：
\begin{equation}
b_k^{\min}\leq b_{ir}\leq B_k,\qquad
b_{jr}\leq b_{ir}+Y_{ir}-e^E_{ij,k}+M(1-x_{ijr}),\quad k\in K^E.
\label{eq:soc}
\end{equation}
同一实体车相邻趟$r\rightarrow r'$满足时间与电量连续：
\begin{align}
T_{d^+_{r'}r'}&\geq T_{d^-_rr}+\sum_{q\in\mathcal Q_{kp}^{rr'}}\Delta_q,\label{eq:trip-time}\\
b_{d^+_{r'}r'}&=b_{d^-_rr}+\sum_{q\in\mathcal Q_{kp}^{rr'}}Y_q.\label{eq:trip-energy}
\end{align}
公共站容量约束为
\begin{equation}
\sum_{k\in K}\sum_{p\in\Omega_k}O_{skpt}z_{kp}\leq C_s,\qquad s\in S,\ t\in T.
\label{eq:station-capacity}
\end{equation}
式\eqref{eq:station-capacity}不施加于车场充电。

设第$m$次重规划时刻为$\tau_m$，待服务客户集合按事件更新：
\begin{equation}
N^{\tau_m}=\left(N^{\tau_{m-1}}\setminus N_{\mathrm{served}}^{\tau_m}\setminus N_{\mathrm{cancel}}^{\tau_m}\right)\cup N_{\mathrm{new}}^{\tau_m}.
\label{eq:dynamic-set}
\end{equation}
已经发生的成本与排放分别累计为
\begin{equation}
\bar C^{\tau_m}=\bar C^{\tau_{m-1}}+\Delta C^{\tau_{m-1}},\qquad
\bar E^{\tau_m}=\bar E^{\tau_{m-1}}+\Delta E^{\tau_{m-1}}.
\label{eq:dynamic-account}
\end{equation}
新一轮优化仅计算剩余问题，已执行前缀作为不可撤销约束。

\section{算法设计}

\subsection{算法框架}

GCI-DMM-VRP同时包含路径、车型、实体车多趟、充电和动态状态继承。本文采用成熟混合遗传搜索作为搜索底座，在其上组织一条面向本问题的串行改进链\cite{ref:vidal-2012,ref:wouda-2024}。底座本身不作为算法创新；问题特化组件的最终贡献由被接受的增量式消融决定。

框架坚持“唯一一本账”：候选是否被接受、能否进入种群、能否存活以及能否被选为父代，均由同一完整评价器给出的完整成本与完整可行性决定。能源消耗、充电、跨趟时间---电量衔接、时变电价和时变碳强度进入同一本账。候选完成失败保留失败原因，不静默冒充初始解。

\begin{figure}[htbp]
  \centering
  \begin{tikzpicture}[
    node distance=5mm and 7mm,
    box/.style={draw=black!55,rounded corners,fill=black!4,align=center,minimum height=8mm,text width=28mm,font=\small},
    arrow/.style={-{Latex[length=2mm]},thick,draw=black!65}
  ]
    \node[box] (snapshot) {静态输入或动态状态快照};
    \node[box,right=of snapshot] (initial) {完整可行起点};
    \node[box,right=of initial] (hgs) {选择、交叉、教育与多样性管理};
    \node[box,below=of hgs] (complete) {实体车多趟与充电补全};
    \node[box,left=of complete] (evaluate) {完整成本与完整可行性评价};
    \node[box,left=of evaluate] (population) {按完整结果入群、生存与父代选择};
    \node[box,below=of evaluate] (audit) {独立验解并保存最好解};
    \draw[arrow] (snapshot)--(initial);
    \draw[arrow] (initial)--(hgs);
    \draw[arrow] (hgs)--(complete);
    \draw[arrow] (complete)--(evaluate);
    \draw[arrow] (evaluate)--(population);
    \draw[arrow] (population.west) -- ++(-6mm,0) |- (hgs.north);
    \draw[arrow] (evaluate)--(audit);
  \end{tikzpicture}
  \caption{算法总体流程。流程图给出已确定的科学合同；\TBD{算法定稿后回填最终组件清单与顺序}。}
  \label{fig:algorithm-flow}
\end{figure}

\subsection{算法步骤}

\textbf{步骤一：构造状态快照与可行起点。} 静态阶段读取订单、车场、有限实体车队和碳强度序列；动态阶段读取式\eqref{eq:dynamic-set}--\eqref{eq:dynamic-account}定义的剩余客户与车辆状态。初始个体先经完整模型检查。

\textbf{步骤二：执行混合遗传搜索主循环。} 算法维护可行与不可行子种群，通过选择、交叉、教育与多样性管理更新候选，罚系数按可行解比例调整。停止条件和预算由本项目收敛记录标定，不直接照搬其他算法的迭代数。

\textbf{步骤三：施加问题特化动作。} 便宜代理只提名有界候选，完整真值逐项评价并决定是否接受。\TBD{算法定稿后列出进入论文的特化动作与出处}。

\textbf{步骤四：完成实体车排班与充电。} 将配送趟分配给具体实体车，按时间顺序衔接多趟并补齐电量、充电会话、分时电价和碳排放。违反时间窗、容量、电量、公共站容量或实体车上限的候选保留失败原因。

\textbf{步骤五：入群、生存与父代选择。} 入群按完整可行性分池，生存选择和父代选择使用同一完整账；全局最好解必须完整可行。

\textbf{步骤六：独立验解与动态调用。} 进入种群的候选重新执行独立检查，只有完整可行且严格改善的候选更新最好解。事件触发后冻结已经执行的路径前缀，更新车辆状态，再求解剩余问题。

\section{数值实验设计与算法有效性}

\subsection{算例、参数与代表解}

当前算例套件保留“China81”历史名称，但权威清单含82个目录条目：成渝27个、京津冀28个、珠三角27个；额外条目是一个京津冀50客户候选。代表实例为京津冀50客户、2车场、97个公共站节点的候选，客户责任按编号等分为两家企业各25个。动态事件流包含10个新增客户事件，接单窗为08:00--10:00，共5个触发批；分时电价和碳强度的信息分辨率为1小时。

\begin{table}[htbp]
  \centering
  \caption{算例规模与结构健康}
  \label{tab:instances}
  \scriptsize
  \begin{tabular}{p{0.30\linewidth}rrrr}
    \toprule
    指标 & 成渝 & 京津冀 & 珠三角 & 全套件\\
    \midrule
    目录条目数（个） & 27 & 28 & 27 & 82\\
    客户规模档（个） & \multicolumn{4}{c}{10、15、20、25、50、75、100、150、200}\\
    每档复本数（个） & 3 & 3（50客户另有候选） & 3 & 3（另有候选）\\
    完成客户数/总客户数 & \TBD{套件正式搜索} & \TBD{套件正式搜索} & \TBD{套件正式搜索} & \TBD{套件正式搜索}\\
    完成需求量/总需求量 & \TBD{套件正式搜索} & \TBD{套件正式搜索} & \TBD{套件正式搜索} & \TBD{套件正式搜索}\\
    每例车场数 & \TBD{汇总} & \TBD{汇总} & \TBD{汇总} & \TBD{汇总}\\
    每例公共站数 & \TBD{汇总} & \TBD{汇总} & \TBD{汇总} & \TBD{汇总}\\
    时间窗宽度 & \TBD{汇总} & \TBD{汇总} & \TBD{汇总} & \TBD{汇总}\\
    上午/下午订单数 & \TBD{汇总} & \TBD{汇总} & \TBD{汇总} & \TBD{汇总}\\
    可争夺结构检查 & \TBD{当前套件复核} & \TBD{当前套件复核} & \TBD{当前套件复核} & \TBD{当前套件复核}\\
    两班次完整服务见证 & \TBD{当前套件复核} & \TBD{当前套件复核} & \TBD{当前套件复核} & \TBD{当前套件复核}\\
    混合车队“两侧”条件 & \TBD{当前套件复核} & \TBD{当前套件复核} & \TBD{当前套件复核} & \TBD{当前套件复核}\\
    \bottomrule
  \end{tabular}
  \tabnote{注：前3行来自当前算例清单。其余行保留已批准表壳，但当前获准数据源没有相应汇总，故不移用旧套件数字。82个目录条目是构造套件条目数，不是企业实测样本数。}
\end{table}

\begin{table}[htbp]
\centering
\caption{关键模型与算法参数及证据身份}
\label{tab:parameters}
\scriptsize
\resizebox{\textwidth}{!}{%
\begin{tabular}{llll}
\toprule
参数 & 数值与单位 & 适用层 & 证据身份/当前处置\\
\midrule
随机种子 & 8 & 代表解复现 & 保存解身份，算法定稿前读数\\
交叉模式 & fast\_only & 交叉 & 当前运行配置；出处待补\\
教育深度 & 1层 & 局部改进 & 标定依据待补，不写成最终值\\
停滞耐心 & 500次无改善 & 停止/重启 & 当前运行配置；待正式标定\\
初始/最小/最大罚系数 & 100.0/0.1/100000.0 & 可行性罚项 & 开源上游默认值\\
罚系数降低/提高倍数 & 0.32/1.34 & 可行性罚项 & 开源上游默认值\\
更新间隔/目标可行比例/容差 & 50/0.43/0.05 & 可行性罚项 & 开源上游默认值\\
修复概率/增强倍数 & 0.8/12 & 不可行解修复 & 开源上游配置\\
最小/最大种群 & 25/65个个体 & 种群管理 & 最大值由25+40得到\\
每代生成数/精英数/邻近数 & 40/4/5 & 种群管理 & 开源上游配置\\
多样性下界/上界 & 0.1/0.5 & 种群多样性 & 开源上游配置\\
锦标赛规模 & 2个个体 & 父代选择 & 当前配置；出处待补\\
充电量/时刻策略 & just\_enough/cost\_plus\_carbon & 充电 & 当前源解实际配置\\
碳价 & 0.07502元/kgCO$_2$e & 总目标 & 全国碳市场价格情景；原始页码待补\\
车场分时电价 & 谷0.56328575、平0.83644275、峰1.14862175元/kWh & 车场充电 & 北京2025-02-12数据版本\\
公共站服务费 & 0.4元/kWh & 公共站充电 & 构造情景，不写成市场实测\\
柴油价 & 7.48元/L & 燃油运营 & 北京地方零售表数据版本\\
逐时碳强度 & 0.15410000--0.64390000 kgCO$_2$e/kWh & 充电排放 & 北京2025-02-12，24个独立小时值\\
燃油/电动车非能源里程成本 & 0.7800/0.9145元/km & 车型与运营成本 & 电动车值含0.2445元/km电池折旧\\
燃油/电动车日固定成本 & 170.00/220.00元/(辆$\cdot$日) & 车辆启用 & 电动车含50元日固定溢价\\
车场单车额定充电功率 & 60.0 kW & 充电可行性 & 中国官方规划取值\\
车辆容积容量 & 7.2 m$^3$ & 路线可行性 & 当前数据集给定；出处待补\\
两车场车队上限 & 3CV+3EV；7CV+7EV & 车型与车辆启用 & 两班次物理见证派生\\
\bottomrule
\end{tabular}
}
\tabnote{注：算法参数只描述表\ref{tab:final-solution}源解的真实运行身份；标为“待补”的来源不因进入本表而升级为正式参数证据。}
\end{table}

\begin{figure}[htbp]
  \centering
  \begin{minipage}{0.48\linewidth}
  \centering
  \begin{tikzpicture}
  \begin{axis}[
    width=\linewidth,height=58mm,
    xlabel={经度},ylabel={纬度},
    tick label style={font=\scriptsize},label style={font=\small},
    legend style={font=\scriptsize,at={(0.02,0.98)},anchor=north west,draw=none,fill=white,fill opacity=.8},
    grid=both,grid style={black!8}]
    \addplot[only marks,mark=+,mark size=1.2pt,gray!45]
      table[x=lon,y=lat,col sep=comma,restrict expr to domain={\thisrow{kind}}{2:2}]{data/figure2_nodes_plot.csv};
    \addlegendentry{公共站}
    \addplot[only marks,mark=o,mark size=1.8pt,blue!65]
      table[x=lon,y=lat,col sep=comma,restrict expr to domain={\thisrow{kind}}{0:0},restrict expr to domain={\thisrow{enterprise}}{0:0}]{data/figure2_nodes_plot.csv};
    \addlegendentry{企业甲客户}
    \addplot[only marks,mark=triangle*,mark size=1.7pt,orange!85!black]
      table[x=lon,y=lat,col sep=comma,restrict expr to domain={\thisrow{kind}}{0:0},restrict expr to domain={\thisrow{enterprise}}{1:1}]{data/figure2_nodes_plot.csv};
    \addlegendentry{企业乙客户}
    \addplot[only marks,mark=square*,mark size=3pt,black]
      table[x=lon,y=lat,col sep=comma,restrict expr to domain={\thisrow{kind}}{1:1}]{data/figure2_nodes_plot.csv};
    \addlegendentry{车场}
  \end{axis}
  \end{tikzpicture}
  \\(a) 节点分布
  \end{minipage}\hfill
  \begin{minipage}{0.48\linewidth}
  \centering
  \begin{tikzpicture}
  \begin{axis}[
    width=\linewidth,height=58mm,
    xlabel={经度},ylabel={纬度},
    tick label style={font=\scriptsize},label style={font=\small},
    grid=both,grid style={black!8},unbounded coords=jump,
    legend style={font=\scriptsize,at={(0.02,0.98)},anchor=north west,draw=none,fill=white,fill opacity=.8}]
    \addplot[thin,blue!60] table[x=lon,y=lat,col sep=comma]{data/figure2_routes_a.csv};
    \addlegendentry{企业甲所属车}
    \addplot[thin,orange!80!black] table[x=lon,y=lat,col sep=comma]{data/figure2_routes_b.csv};
    \addlegendentry{企业乙所属车}
    \addplot[only marks,mark=o,mark size=1.3pt,black!65]
      table[x=lon,y=lat,col sep=comma,restrict expr to domain={\thisrow{kind}}{0:0}]{data/figure2_nodes_plot.csv};
    \addplot[only marks,mark=square*,mark size=3pt,black]
      table[x=lon,y=lat,col sep=comma,restrict expr to domain={\thisrow{kind}}{1:1}]{data/figure2_nodes_plot.csv};
  \end{axis}
  \end{tikzpicture}
  \\(b) 路线及车型
  \end{minipage}
  \caption{代表实例的节点分布与路线。节点表含2个车场、50个客户和97个公共站节点；保存解包含16趟路线，均由燃油车执行。路线未作手工平滑或美化。}
  \label{fig:nodes-routes}
\end{figure}

\begin{table}[htbp]
  \centering
  \caption{代表实例解的成本、排放与资源分解}
  \label{tab:final-solution}
  \small
  \begin{tabular}{llr}
    \toprule
    类别 & 指标（单位） & 保存解数值\\
    \midrule
    服务 & 完成客户数/总客户数（个） & 50/50\\
    服务 & 完成需求量/总需求量（kg） & 13264/13264\\
    资源 & 启用燃油车/电动车（辆） & 8/0\\
    资源 & 启用实体车总数（辆） & 8\\
    资源 & 配送趟次（趟） & 16\\
    里程 & 总里程（km） & 935.530\\
    成本 & 车辆固定成本（元） & 1360.000\\
    成本 & 非能源里程成本（元） & 729.713\\
    成本 & 燃油成本（元） & 919.139\\
    成本 & 电费（元） & 0.000\\
    成本 & 碳价折算成本（元） & 24.354\\
    成本 & 单目标总账（元） & 3033.206\\
    能源 & 燃油消耗（L） & 122.879\\
    能源 & 电能消耗（kWh） & 0.000\\
    排放 & 燃油车直接排放（kgCO$_2$e） & 324.636\\
    排放 & 电动车充电排放（kgCO$_2$e） & 0.000\\
    可行性 & 完整评价可行 & 是\\
    可行性 & 硬约束违规数（项） & 0\\
    \bottomrule
  \end{tabular}
  \tabnote{注：保存解来自协同效率批的联合运行，公平约束关闭。该解用于展示一份可执行方案如何闭合服务、资源、成本与排放，不承担算法优越性证明。总里程由8辆实体车的逐车里程求和得到，与联合解分解字段一致。}
\end{table}

\begin{table}[htbp]
  \centering
  \caption{代表解实体车辆的多趟安排}
  \label{tab:vehicle-schedules}
  \scriptsize
  \resizebox{\textwidth}{!}{%
  \begin{tabular}{llllrr}
    \toprule
    实体车 & 所属车场 & 车型 & 趟数 & 服务客户数 & 里程（km）\\
    \midrule
    CV\_A\_1 & 车场A & 燃油车 & 2 & 6 & 151.083\\
    CV\_A\_2 & 车场A & 燃油车 & 1 & 2 & 83.613\\
    CV\_B\_1 & 车场B & 燃油车 & 1 & 4 & 56.999\\
    CV\_B\_2 & 车场B & 燃油车 & 3 & 10 & 151.064\\
    CV\_B\_3 & 车场B & 燃油车 & 3 & 8 & 162.811\\
    CV\_B\_4 & 车场B & 燃油车 & 2 & 6 & 93.774\\
    CV\_B\_5 & 车场B & 燃油车 & 2 & 7 & 119.043\\
    CV\_B\_6 & 车场B & 燃油车 & 2 & 7 & 117.142\\
    \bottomrule
  \end{tabular}
  }
  \tabnote{注：逐车客户序列已经保存，但当前展品数据没有逐趟发车/返回时刻、油耗、直接排放与运营成本字段；这些项目均为\TBD{重新完整评价保存解后提取}。}
\end{table}


\subsection{公开算例上的求解质量}

公开实验采用28个多车场带时间窗算例检验基础路径搜索能力。当前每题只运行一次，单题时限约120秒，预算未按收敛曲线完成标定。因此本节只展示逐题解、服务完整性和原始精度路线复核，不写“优于”“领先”或“胜出”。公开最好值及其版本尚未补齐。

\begin{table}[htbp]
  \centering
  \caption{公开28题的单次求解结果与服务完整性}
  \label{tab:public28}
  \scriptsize
  \resizebox{\textwidth}{!}{%
  \begin{tabular}{lrrrrrlc}
    \toprule
    算例 & 本文Best & 原始精度距离 & 客户完成 & 路线可行 & 圈数 & 运行时间(s) & 公开最好值/版本\\
    \midrule
    PR11A & 6776.761 & 6776.756615 & 360/360 & 31/31 & 13679 & 120.004 & \TBD{BKS}\\
    PR11B & 4820.376 & 4820.369984 & 360/360 & 25/25 & 13680 & 115.783 & \TBD{BKS}\\
    PR12A & 8251.671 & 8251.667034 & 480/480 & 41/41 & 10006 & 120.013 & \TBD{BKS}\\
    PR12B & 6041.261 & 6041.261577 & 480/480 & 33/33 & 11471 & 120.013 & \TBD{BKS}\\
    PR13A & 9722.412 & 9722.403762 & 600/600 & 48/48 & 6509 & 120.009 & \TBD{BKS}\\
    PR13B & 7293.739 & 7293.739728 & 600/600 & 41/41 & 7644 & 120.009 & \TBD{BKS}\\
    PR14A & 11280.749 & 11280.746232 & 720/720 & 57/57 & 4630 & 120.007 & \TBD{BKS}\\
    PR14B & 8825.714 & 8825.706870 & 720/720 & 53/53 & 5646 & 120.008 & \TBD{BKS}\\
    PR15A & 13229.546 & 13229.527282 & 840/840 & 69/69 & 3034 & 120.026 & \TBD{BKS}\\
    PR15B & 10522.498 & 10522.511209 & 840/840 & 64/64 & 3456 & 120.017 & \TBD{BKS}\\
    PR16A & 14633.726 & 14633.711735 & 960/960 & 78/78 & 2571 & 120.022 & \TBD{BKS}\\
    PR16B & 11762.652 & 11762.646734 & 960/960 & 71/71 & 2701 & 120.020 & \TBD{BKS}\\
    PR17A & 6316.228 & 6316.224494 & 360/360 & 30/30 & 11868 & 77.273 & \TBD{BKS}\\
    PR17B & 4796.924 & 4796.935613 & 360/360 & 26/26 & 17694 & 120.002 & \TBD{BKS}\\
    PR18A & 8365.857 & 8365.858590 & 520/520 & 44/44 & 8884 & 120.018 & \TBD{BKS}\\
    PR18B & 6512.750 & 6512.754096 & 520/520 & 38/38 & 9291 & 120.002 & \TBD{BKS}\\
    PR19A & 10831.511 & 10831.516882 & 700/700 & 57/57 & 3311 & 120.005 & \TBD{BKS}\\
    PR19B & 8199.189 & 8199.192843 & 700/700 & 52/52 & 3996 & 120.017 & \TBD{BKS}\\
    PR20A & 12204.194 & 12204.201413 & 880/880 & 70/70 & 2668 & 120.002 & \TBD{BKS}\\
    PR20B & 10324.916 & 10324.931913 & 880/880 & 68/68 & 2842 & 120.010 & \TBD{BKS}\\
    PR21A & 6251.424 & 6251.429902 & 420/420 & 35/35 & 11271 & 120.007 & \TBD{BKS}\\
    PR21B & 4893.444 & 4893.457803 & 420/420 & 30/30 & 12204 & 120.002 & \TBD{BKS}\\
    PR22A & 8121.877 & 8121.888102 & 600/600 & 46/46 & 4903 & 120.010 & \TBD{BKS}\\
    PR22B & 6516.749 & 6516.762328 & 600/600 & 42/42 & 5440 & 120.015 & \TBD{BKS}\\
    PR23A & 10138.318 & 10138.324009 & 780/780 & 62/62 & 3467 & 120.014 & \TBD{BKS}\\
    PR23B & 8480.755 & 8480.769299 & 780/780 & 58/58 & 4080 & 120.002 & \TBD{BKS}\\
    PR24A & 12149.427 & 12149.435442 & 960/960 & 81/81 & 2301 & 120.016 & \TBD{BKS}\\
    PR24B & 10905.516 & 10905.518704 & 960/960 & 75/75 & 2446 & 120.009 & \TBD{BKS}\\
    \bottomrule
  \end{tabular}
  }
  \tabnote{注：28题均完成全部客户，全部路线通过原始精度复核，硬约束违规为零。每题只有一次运行，约120秒；预算未按收敛数据标定。“本文Best”和“原始精度距离”的显示差异来自记录精度，不作算法比较。公开最好值与版本尚未补齐，故不计算差距。}
\end{table}

\subsection{算法组件贡献分析}

增量式消融原设计从路线基准开始，逐步加入跨场协同、实体车多趟、整车换型与时变碳充电。首批15个单元中9个成功，但两个基准配置各自0/3成功，整包判定为未接受。由于基准缺失，现有终值不能形成相对基准的算法贡献链，表\ref{tab:ablation}全部保留占位。

\begin{table}[htbp]
  \centering
  \caption{串行改进链的增量式消融}
  \label{tab:ablation}
  \scriptsize
  \resizebox{\textwidth}{!}{%
  \begin{tabular}{lp{0.18\linewidth}rrrrrr}
    \toprule
    算法臂 & 累加组件 & Best & Avg & Std & 相对上一臂 & 相对基准 & 完整可行运行数\\
    \midrule
    基准 & 路线基准 & \TBD{重做} & \TBD{重做} & \TBD{重做} & --- & --- & \TBD{重做}\\
    累加一 & 跨场协同 & \TBD{重做} & \TBD{重做} & \TBD{重做} & \TBD{重做} & \TBD{重做} & \TBD{重做}\\
    累加二 & 实体车多趟 & \TBD{重做} & \TBD{重做} & \TBD{重做} & \TBD{重做} & \TBD{重做} & \TBD{重做}\\
    累加三 & 整车换型 & \TBD{重做} & \TBD{重做} & \TBD{重做} & \TBD{重做} & \TBD{重做} & \TBD{重做}\\
    累加四 & 时变碳充电 & \TBD{重设计} & \TBD{重设计} & \TBD{重设计} & \TBD{重设计} & \TBD{重设计} & \TBD{重设计}\\
    \bottomrule
  \end{tabular}
  }
  \tabnote{注：首批消融整包accepted=false；两个基准臂均0/3成功，不能充当有效分母。时变碳充电属于机制实验，是否保留在算法增量链中仍需重设计。本表不搬入首批终值。}
\end{table}

\begin{figure}[htbp]
  \centering
  \begin{tikzpicture}
  \begin{axis}[
    width=0.92\textwidth,height=72mm,
    xlabel={完整模型评价次数},ylabel={当前最好完整目标值（元）},
    grid=both,grid style={black!8},
    tick label style={font=\scriptsize},label style={font=\small},
    legend columns=3,legend style={font=\scriptsize,at={(0.5,1.02)},anchor=south,draw=none}]
    \addplot[blue!45,thin] table[x=evaluations,y=cost,col sep=comma,restrict expr to domain={\thisrow{armcode}}{2:2},restrict expr to domain={\thisrow{seed}}{1:1}]{data/convergence_plot.csv};
    \addlegendentry{实体车多趟}
    \addplot[blue!65,dashed,thin] table[x=evaluations,y=cost,col sep=comma,restrict expr to domain={\thisrow{armcode}}{2:2},restrict expr to domain={\thisrow{seed}}{2:2}]{data/convergence_plot.csv};
    \addplot[blue!85,dotted,thin] table[x=evaluations,y=cost,col sep=comma,restrict expr to domain={\thisrow{armcode}}{2:2},restrict expr to domain={\thisrow{seed}}{3:3}]{data/convergence_plot.csv};
    \addplot[orange!55!black,thin] table[x=evaluations,y=cost,col sep=comma,restrict expr to domain={\thisrow{armcode}}{3:3},restrict expr to domain={\thisrow{seed}}{1:1}]{data/convergence_plot.csv};
    \addlegendentry{整车换型}
    \addplot[orange!75!black,dashed,thin] table[x=evaluations,y=cost,col sep=comma,restrict expr to domain={\thisrow{armcode}}{3:3},restrict expr to domain={\thisrow{seed}}{2:2}]{data/convergence_plot.csv};
    \addplot[orange!95!black,dotted,thin] table[x=evaluations,y=cost,col sep=comma,restrict expr to domain={\thisrow{armcode}}{3:3},restrict expr to domain={\thisrow{seed}}{3:3}]{data/convergence_plot.csv};
    \addplot[green!45!black,thin] table[x=evaluations,y=cost,col sep=comma,restrict expr to domain={\thisrow{armcode}}{4:4},restrict expr to domain={\thisrow{seed}}{1:1}]{data/convergence_plot.csv};
    \addlegendentry{时变碳充电}
    \addplot[green!65!black,dashed,thin] table[x=evaluations,y=cost,col sep=comma,restrict expr to domain={\thisrow{armcode}}{4:4},restrict expr to domain={\thisrow{seed}}{2:2}]{data/convergence_plot.csv};
    \addplot[green!85!black,dotted,thin] table[x=evaluations,y=cost,col sep=comma,restrict expr to domain={\thisrow{armcode}}{4:4},restrict expr to domain={\thisrow{seed}}{3:3}]{data/convergence_plot.csv};
  \end{axis}
  \end{tikzpicture}
  \caption*{收敛曲线工作图（非正式编号）。首批消融中三个可行配置各含三个随机种子的原始轨迹，共使用72行保存数据。路线基准与跨场协同配置均未产生可行曲线，故缺席；整包未被接受，本图只展示已有搜索过程，不支撑组件优劣结论。}
\end{figure}

\section{时变碳强度下的减排杠杆}

本章按充电时刻、车型指派和动态发车三个层级组织。三节共用式\eqref{eq:system-emission}和式\eqref{eq:objective}。当前充电时刻层已有固定排班作用空间数据；车型层只有两档成功探针；动态层没有可用结果，因此各节的结论强度不同。

\subsection{充电时刻层}

\subsubsection{这一层要回答什么}

电动配送车的排放不发生在路上，而发生在充电的那一刻：同样的一度电，在电网清洁时段充和在高碳时段充，间接排放可以相差数倍。因此，在路线、车型与充电量都不变的情况下，充电时刻本身就是一个减排杠杆。

这一层要回答三件事：第一，这个杠杆在本文算例中有多大作用空间；第二，若使用日均排放因子核算，这个作用空间是否还看得见；第三，在现行碳价下，动用这个杠杆是否划算。

\subsubsection{实验设计}

四个对照臂固定同一实例、同一批订单、同一套排班与路线、同一车型指派、同一充电地点、同一总充电量、同一可行充电窗口、同一车辆初始状态与服务要求，只改变碳强度的表示方式与充电选择规则：固定平均碳强度使用全天24小时等权均值0.4393~kgCO$_2$e/kWh并沿用有空即充排程；另三臂均使用逐小时值，分别采取有空即充、电费最省和碳感知规则。

算例为京津冀代表实例（\texttt{cn-jjj-50c-01-DEPOTSEARCH-d996f755bd}），电网日为2025-02-12，代表城市为北京。碳强度取北京省级逐小时投影值，在0.1541--0.6439~kgCO$_2$e/kWh之间变动；最清洁时段为13:00，最肮脏时段为01:00。全天总充电量为189.279~kWh，两条已有时变策略的路线、车型、客户服务关系、充电地点和每次充电量均相同。

\subsubsection{时变电网碳强度与充电选择规则对比}

\begin{table}[htbp]
  \centering
  \caption{时变电网碳强度与充电选择规则对比}
  \label{tab:carbon-charging}
  \scriptsize
  \resizebox{\textwidth}{!}{%
  \begin{tabular}{lrrrrr}
    \toprule
    指标（单位） & 固定平均碳强度 & 时变---有空即充 & 时变---电费最省 & 时变---碳感知 & 本文相对固定基准\\
    \midrule
    完成客户数（个） & 50/50 & 50/50 & \TBD{该策略} & 50/50 & 0\\
    完成需求量（kg） & \TBD{源表缺字段} & 同左 & \TBD{该策略} & 同左 & 0\\
    可移动充电量（kWh） & 0.000 & 0.000 & \TBD{该策略} & 174.732 & ---\\
    充电电费（元） & 141.525 & 141.525 & \TBD{该策略} & 168.049 & +18.742\%\\
    碳价折算成本（元，碳价0.075） & 6.236 & 6.679 & \TBD{该策略} & 2.247 & $-63.961\%$\\
    充电侧总账（元） & 147.761 & 148.203 & \TBD{该策略} & 170.297 & +15.252\%\\
    全系统总账（元） & 3032.527 & 3032.970 & \TBD{该策略} & 3055.063 & +0.743\%\\
    充电排放（kgCO$_2$e） & 83.150 & 89.048 & \TBD{该策略} & 29.967 & $-63.961\%$\\
    配送作业阶段总排放（kgCO$_2$e） & 183.917 & 189.816 & \TBD{该策略} & 130.734 & $-28.917\%$\\
    \bottomrule
  \end{tabular}
  }
  \tabnote{注：固定平均碳强度列由总充电量189.27888317792736~kWh乘日均因子0.4393得到；其充电排程与有空即充相同。四列应保持相同服务量；电费最省列尚未产生。充电侧总账等于电费与充电排放碳成本之和；全系统总账另含固定不变的其他成本2877.208798592812元。固定排班数据来自探针源解上的精确枚举，展示作用空间，不冒充多种子正式统计。}
\end{table}

\subsubsection{读表}

第一，作用空间是真实存在的，而且很大。在排班完全不变的前提下，全天189.279~kWh中有174.732~kWh（92.315\%）可以改变充电时刻；把它们从当前时段移到合法窗口内的最清洁组合，充电侧排放由89.048降至29.967~kgCO$_2$e，降幅66.348\%；由于燃油车直接排放不受影响，折算到配送作业阶段总排放为下降28.917\%。

第二，日均口径会让这个作用空间在账面上完全消失。在固定平均碳强度下，无论何时充电，充电排放的账面值恒为83.150~kgCO$_2$e；用日均排放因子核算，上述59.082~kgCO$_2$e差异无法呈现，同时还把当前排班的逐时复算排放低估5.898~kgCO$_2$e，即6.623\%。

第三，在现行碳价下，动用这个杠杆并不划算。最清洁时段不是最便宜时段：碳感知排程使电费上升26.525元，而节省的排放按0.075元/kg碳价仅折算4.431元，完整目标净增22.094元。使两者持平所需的配对碳价为0.4489498210731408元/kg，正文显示为0.449元/kg。

碳价扫描给出这条判断的完整形状：在0.075、0.20、0.60、1.10、1.50和2.10元/kg六个点，碳感知相对有空即充的完整目标变化依次为+22.094、+14.708、$-8.924$、$-38.465$、$-62.098$和$-97.547$元。

\subsubsection{时变碳强度与充电响应}

\begin{figure}[htbp]
  \centering
  \begin{minipage}{0.32\linewidth}
  \centering
  \begin{tikzpicture}
    \begin{axis}[
      width=\linewidth,height=47mm,xmin=0,xmax=23,
      xlabel={小时},ylabel={电价（元/kWh）},
      axis y line*=left,axis x line*=bottom,
      tick label style={font=\tiny},label style={font=\scriptsize},
      xtick={0,6,12,18,23},grid=major,grid style={black!8}]
      \addplot[orange!85!black,thick,const plot mark left] table[x=hour,y=price,col sep=comma]{data/hourly_plot.csv};
    \end{axis}
    \begin{axis}[
      width=\linewidth,height=47mm,xmin=0,xmax=23,
      axis y line*=right,axis x line=none,
      ylabel={碳强度（kgCO$_2$e/kWh）},
      tick label style={font=\tiny},label style={font=\scriptsize}]
      \addplot[blue!70,thick,const plot mark left] table[x=hour,y=carbon,col sep=comma]{data/hourly_plot.csv};
    \end{axis}
  \end{tikzpicture}
  \\(a) 错位
  \end{minipage}\hfill
  \begin{minipage}{0.32\linewidth}
  \centering
  \fbox{\parbox[c][43mm][c]{0.91\linewidth}{\centering\TBD{逐车发车/返回时刻与充电窗口数据}}}
  \\(b) 窗口
  \end{minipage}\hfill
  \begin{minipage}{0.32\linewidth}
  \centering
  \begin{tikzpicture}
    \begin{axis}[
      width=\linewidth,height=47mm,
      xlabel={碳价（元/kg）},ylabel={完整目标变化（元）},
      tick label style={font=\tiny},label style={font=\scriptsize},
      grid=major,grid style={black!8},
      xmin=0,xmax=2.15]
      \addplot[green!55!black,thick,const plot mark right,mark=*] table[x=price,y=delta,col sep=comma]{data/carbon_price_scan.csv};
      \addplot[black!55,dashed] coordinates {(0.4489498210731408,-105) (0.4489498210731408,30)};
      \addplot[black!40] coordinates {(0,0) (2.15,0)};
      \node[font=\tiny,anchor=south west] at (axis cs:0.4489498210731408,0) {0.449};
    \end{axis}
  \end{tikzpicture}
  \\(c) 代价
  \end{minipage}
  \caption{时变电网碳强度下的充电响应。（a）场内分时电价与电网碳强度均按24个逐小时值画阶梯；（b）批准形态为逐车时间线，但当前展品数据缺少逐车发车与返回时刻，故保留同尺寸占位；（c）碳感知相对有空即充的完整目标变化，翻转点为0.449元/kg。}
  \label{fig:carbon-charging}
\end{figure}

图\ref{fig:carbon-charging}(a)显示分时电价把充电往深夜赶，而深夜接近电网最脏时段；全天最清洁的13:00落在平段电价上，低价与低碳不在同一时刻。(b)按批准形态保留逐车时间线，但本轮数据缺少逐车发车与返回时刻，故不把聚合充电动作伪装成完整时间线；已知10次充电中9次改变时刻，174.732~kWh被移动，另有14.547~kWh因窗口限制保持原时刻。(c)显示完整目标变化随碳价由正转负，配对算式给出的翻转点为0.449元/kg。因此本层的结论不是“挪不动”，而是“挪得动、但现行碳价下不值得挪”。

\subsubsection{小结}

在固定排班下，充电时刻层的减排作用空间为59.082~kgCO$_2$e/日，占配送作业阶段总排放的28.917\%；日均排放因子核算会使该空间在账面上不可见；动用该空间的盈亏平衡碳价为0.449元/kg，约为0.075元/kg的6倍。这说明信息分辨率与价格水平是两个独立前提：只提高碳核算的时间分辨率而不改变价格信号，减排潜力可见但不会被执行。

\subsection{车型指派层}

车型层已登记四个近端点车队档位：18/2、13/7、7/13和2/18（燃油车/电动车）。参数类要求两种车型各至少保留一辆，因此真正的纯燃油或纯电端点当前不可表示。前两档完成探针运行，后两档因未登记实体车辆而失败。

\begin{table}[htbp]
  \centering
  \caption{车队可用结构与实际车型指派}
  \label{tab:fleet-levels}
  \scriptsize
  \resizebox{\textwidth}{!}{%
  \begin{tabular}{lrrrr}
    \toprule
    指标（单位） & 18CV/2EV & 13CV/7EV & 7CV/13EV & 2CV/18EV\\
    \midrule
    完成客户数（个） & 50/50 & 50/50 & \TBD{搜索未完成} & \TBD{搜索未完成}\\
    完成需求量（kg） & 13264/13264 & 13264/13264 & \TBD{搜索未完成} & \TBD{搜索未完成}\\
    实际派遣CV/EV（辆） & 9/0 & 10/0 & \TBD{搜索未完成} & \TBD{搜索未完成}\\
    配送趟次（趟） & 17 & 16 & \TBD{搜索未完成} & \TBD{搜索未完成}\\
    总里程（km） & 1013.128 & 936.115 & \TBD{搜索未完成} & \TBD{搜索未完成}\\
    车辆固定成本（元） & 1530.000 & 1700.000 & \TBD{搜索未完成} & \TBD{搜索未完成}\\
    非能源里程成本（元） & 790.240 & 730.169 & \TBD{搜索未完成} & \TBD{搜索未完成}\\
    燃油成本（元） & 992.863 & 918.361 & \TBD{搜索未完成} & \TBD{搜索未完成}\\
    电费（元） & 0.000 & 0.000 & \TBD{搜索未完成} & \TBD{搜索未完成}\\
    碳价折算成本（元） & 26.308 & 24.334 & \TBD{搜索未完成} & \TBD{搜索未完成}\\
    单目标总账（元） & 3339.410 & 3372.864 & \TBD{搜索未完成} & \TBD{搜索未完成}\\
    燃油车直接排放（kgCO$_2$e） & 350.675 & 324.361 & \TBD{搜索未完成} & \TBD{搜索未完成}\\
    电动车充电排放（kgCO$_2$e） & 0.000 & 0.000 & \TBD{搜索未完成} & \TBD{搜索未完成}\\
    运行身份/状态 & 探针完成 & 探针完成 & \TBD{未登记实体车} & \TBD{未登记实体车}\\
    \bottomrule
  \end{tabular}
  }
  \tabnote{注：前两档的verdict为PROBE\_RUN\_COMPLETE，并非正式多种子结果；两档实际派遣电动车均为0。后两档失败的近端原因是完成解引用了未登记的燃油实体车，因此写“搜索未完成”，不写成“模型不可行”。}
\end{table}

现有两档只说明在这两次探针中，增加可用电动车名额没有转化为实际电动车派遣。由于电动车较多的两档未完成，当前不能据此写出完整车型构成规律。

\begin{figure}[htbp]
  \centering
  \fbox{\parbox[c][58mm][c]{0.86\linewidth}{\centering
    \TBD{混合车队碳价扫描：横轴碳价，纵轴燃油车/电动车实际派遣数；需完整碳价序列下的可比正式结果}}}
  \caption{碳价变化下的车型指派。批准形态为单幅图，分别画燃油车和电动车数量；现有四档探针不是碳价扫描，故不把它们拼成该图。}
  \label{fig:fleet-carbon-price}
\end{figure}

\subsection{动态发车层}

动态实验设计要求所有策略面对同一条不泄漏事件流、同一初始计划、车辆状态、外包规则、计算预算和服务要求；已经执行的路线前缀冻结。完全信息静态只作理论参照，机械在线与碳感知在线才构成同信息条件下的比较。现有动态结果包accepted=false，求解器没有进入，因而本节不写任何结果句。

\begin{table}[htbp]
  \centering
  \caption{动态需求响应方式对比}
  \label{tab:dynamic}
  \scriptsize
  \resizebox{\textwidth}{!}{%
  \begin{tabular}{lrrrr}
    \toprule
    指标（单位） & 完全信息静态 & 机械在线 & 碳感知在线 & 本文相对机械在线\\
    \midrule
    响应订单数（个） & \TBD{动态正式批} & \TBD{动态正式批} & \TBD{动态正式批} & \TBD{比较}\\
    完成客户数（个） & \TBD{动态正式批} & \TBD{动态正式批} & \TBD{动态正式批} & \TBD{比较}\\
    完成需求量（kg） & \TBD{动态正式批} & \TBD{动态正式批} & \TBD{动态正式批} & \TBD{比较}\\
    拒绝/外包订单数（个） & \TBD{动态正式批} & \TBD{动态正式批} & \TBD{动态正式批} & \TBD{比较}\\
    外包成本（元） & \TBD{动态正式批} & \TBD{动态正式批} & \TBD{动态正式批} & \TBD{比较}\\
    实际车辆数（辆） & \TBD{动态正式批} & \TBD{动态正式批} & \TBD{动态正式批} & \TBD{比较}\\
    重规划次数（次） & \TBD{动态正式批} & \TBD{动态正式批} & \TBD{动态正式批} & \TBD{比较}\\
    总里程（km） & \TBD{动态正式批} & \TBD{动态正式批} & \TBD{动态正式批} & \TBD{比较}\\
    单目标总账（元） & \TBD{动态正式批} & \TBD{动态正式批} & \TBD{动态正式批} & \TBD{比较}\\
    燃油车直接排放（kgCO$_2$e） & \TBD{动态正式批} & \TBD{动态正式批} & \TBD{动态正式批} & \TBD{比较}\\
    电动车充电排放（kgCO$_2$e） & \TBD{动态正式批} & \TBD{动态正式批} & \TBD{动态正式批} & \TBD{比较}\\
    实际运行时间（s） & \TBD{动态正式批} & \TBD{动态正式批} & \TBD{动态正式批} & \TBD{比较}\\
    \bottomrule
  \end{tabular}
  }
  \tabnote{注：动态包accepted=false，求解器未进入；本表不使用该包声称路线变化、成本差或排放差。完全信息静态预先知道全日订单，只作理论参照。}
\end{table}

\DraftNote{动态状态展品待填字段：事件时刻与类型、触发时刻、待服务客户、车辆位置、剩余容量、剩余电量、最晚服务余量、冻结前缀和新增车辆。以上字段必须来自同一次被接受的动态运行，不能从当前失败包拼接。}

\begin{figure}[htbp]
  \centering
  \begin{minipage}{0.46\linewidth}
    \centering\fbox{\parbox[c][48mm][c]{0.91\linewidth}{\centering\TBD{触发前路线、车辆状态与冻结弧}}}\\(a) 触发前
  \end{minipage}\hfill
  \begin{minipage}{0.46\linewidth}
    \centering\fbox{\parbox[c][48mm][c]{0.91\linewidth}{\centering\TBD{触发后路线、虚拟起点与重规划弧}}}\\(b) 触发后
  \end{minipage}
  \caption{代表动态事件触发前后的路线。两格保留统一坐标和尺寸；当前没有可用动态结果，故不画路线。}
  \label{fig:dynamic-routes}
\end{figure}

综合本章，充电时刻层已有固定排班作用空间数据；车型层只有两个成功探针且均未派遣电动车；动态层没有结果。因而“主要减排杠杆是否随条件转移”的全文结论仍为\TBD{三层可比正式结果}。

\section{协同效率与收益分配}

协同和公平使用同一代表实例，但回答不同问题。协同是效率问题：关闭公平约束，比较两家各自优化后的单干成本之和与联合优化成本。公平是分配问题：固定已经得到的联合方案，只改变联合总成本在两家之间的承担额，不重新求路线。

\subsection{多车场协同效率}

两家企业的最低单干成本分别取各自保存批次中的最好值，合计3922.0786726680826元；关闭公平约束的联合最好值为3033.2059767107694元。两边均完成50/50客户和13264/13264~kg需求，物理违规为零。

\begin{table}[htbp]
  \centering
  \caption{两家单干之和与联合协同效率对比}
  \label{tab:synergy}
  \small
  \begin{tabular}{lrrr}
    \toprule
    指标（单位） & 两家单干之和 & 联合协同 & 联合相对单干\\
    \midrule
    完成客户数（个） & 50/50 & 50/50 & 0\\
    完成需求量（kg） & 13264/13264 & 13264/13264 & 0\\
    启用CV/EV（辆） & 2/6 & 8/0 & 构成变化\\
    启用实体车数（辆） & 8 & 8 & 0\\
    配送趟次（趟） & 16 & 16 & 0\\
    总里程（km） & 1717.715 & 935.530 & $-45.536\%$\\
    运营成本（元） & 3897.623 & 3008.852 & $-22.803\%$\\
    碳价折算成本（元） & 24.456 & 24.354 & $-0.417\%$\\
    单目标总账（元） & 3922.079 & 3033.206 & $-22.663\%$\\
    燃油车直接排放（kgCO$_2$e） & 136.666 & 324.636 & +137.540\%\\
    电动车充电排放（kgCO$_2$e） & 189.328 & 0.000 & $-100.000\%$\\
    配送作业阶段总排放（kgCO$_2$e） & 325.994 & 324.636 & $-0.417\%$\\
    \bottomrule
  \end{tabular}
  \tabnote{注：单干行由企业甲和企业乙各自保存最好解相加；联合行取关闭公平约束的联合最好解。两者服务量相同。联合方案节约888.8726959573132元，即22.663306122634136\%。车型构成同时变化，故本表是整体协同效率比较，不把成本变化单独归因于跨场路线。}
\end{table}

\subsection{协同收益的事后分配}

联合方案形成888.8726959573132元节约。路线归属记账下，企业甲不满足参与条件，企业乙满足；事后分摊不改变联合总成本、路线、服务量和排放，只重新划分承担额。分摊后两家各节约444.436347978656元，参与条件由1/2变为2/2满足。

\begin{table}[htbp]
  \centering
  \caption{有无公平分配的企业级成本承担对比}
  \label{tab:allocation}
  \small
  \begin{tabular}{lrrr}
    \toprule
    分摊方法/指标 & 企业甲 & 企业乙 & 联盟合计\\
    \midrule
    单干成本（元） & 2068.897 & 1853.181 & 3922.079\\
    无公平分配：路线归属成本（元） & 754.041 & 2279.165 & 3033.206\\
    公平分配后成本（元） & 1624.461 & 1408.745 & 3033.206\\
    分配后节约量（元） & 444.436 & 444.436 & 888.873\\
    分配后节约百分比 & 21.482\% & 23.982\% & 22.663\%\\
    路线归属记账下参与条件 & 否 & 是 & 1/2满足\\
    分配后参与条件 & 是 & 是 & 2/2满足\\
    Shapley与核仁规则间极差（元） & 0.000 & 0.000 & 0.000\\
    分配和校验误差（元） & \multicolumn{3}{c}{0.000（显示精度）}\\
    \bottomrule
  \end{tabular}
  \tabnote{注：公平分配发生在同一联合最优解之上，不重新求路线。两方成本博弈中，本数据下Shapley分摊与核仁分摊相同，均落在转归集内；这不是对多方规则差异的证明。旧探索包中的943.967428元和471.983714元不用于本表。}
\end{table}

\section{补充机制与边界}

\subsection{实体车辆多趟}

多趟机制用于区分“配送趟数”和“启用实体车数”，并保证相邻趟的返回、补能和再发车连续。当前代表解展示了8辆实体车执行16趟任务，但没有符合批准合同的单趟/多趟正式对照，因此表\ref{tab:multitrip}只保留结构。

\begin{table}[htbp]
  \centering
  \caption{实体车辆单趟与多趟执行对比}
  \label{tab:multitrip}
  \small
  \begin{tabular}{lrrr}
    \toprule
    指标（单位） & 单趟实体车 & 实体车多趟 & 多趟相对单趟\\
    \midrule
    完成客户数（个） & \TBD{正式两臂} & \TBD{正式两臂} & \TBD{比较}\\
    完成需求量（kg） & \TBD{正式两臂} & \TBD{正式两臂} & \TBD{比较}\\
    启用实体车数（辆） & \TBD{正式两臂} & \TBD{正式两臂} & \TBD{比较}\\
    配送趟次（趟） & \TBD{正式两臂} & \TBD{正式两臂} & \TBD{比较}\\
    平均趟/车（趟/辆） & \TBD{正式两臂} & \TBD{正式两臂} & \TBD{比较}\\
    车辆固定成本（元） & \TBD{正式两臂} & \TBD{正式两臂} & \TBD{比较}\\
    总里程（km） & \TBD{正式两臂} & \TBD{正式两臂} & \TBD{比较}\\
    充电量（kWh） & \TBD{正式两臂} & \TBD{正式两臂} & \TBD{比较}\\
    单目标总账（元） & \TBD{正式两臂} & \TBD{正式两臂} & \TBD{比较}\\
    衔接可行状态 & \TBD{逐车复核} & \TBD{逐车复核} & \TBD{比较}\\
    \bottomrule
  \end{tabular}
  \tabnote{注：现有代表解只证明多趟结构能够出现，不证明多趟相对单趟的机制增量。正式对照必须保持订单、车型、运营时域、预算和服务要求相同。}
\end{table}

\subsection{非线性充电}

非线性充电在本文中的正面叙事不是“必须进入高SOC区间”，而是算法是否利用边际充电功率下降的信息，在后段耗时长、收益小的区间之前停止。若建模了非线性曲线却仍盲目充满，反而说明算法没有利用该信息。当前获准数据源没有完整线性影子/分段/非线性回算表，因此不把历史读数搬入。

\begin{table}[htbp]
  \centering
  \caption{不同充电功率函数的账面求解与非线性回算对比}
  \label{tab:nonlinear-charging}
  \scriptsize
  \resizebox{\textwidth}{!}{%
  \begin{tabular}{lrrrrr}
    \toprule
    指标（单位） & 恒功率S & 恒功率R & 分段SOC功率S & 分段SOC功率R & 非线性SOC功率R\\
    \midrule
    完成客户数（个） & \TBD{正式S/R} & \TBD{正式S/R} & \TBD{正式S/R} & \TBD{正式S/R} & \TBD{正式S/R}\\
    完成需求量（kg） & \TBD{正式S/R} & \TBD{正式S/R} & \TBD{正式S/R} & \TBD{正式S/R} & \TBD{正式S/R}\\
    可行运行数（次） & \TBD{正式S/R} & \TBD{正式S/R} & \TBD{正式S/R} & \TBD{正式S/R} & \TBD{正式S/R}\\
    充电会话数（次） & \TBD{正式S/R} & \TBD{正式S/R} & \TBD{正式S/R} & \TBD{正式S/R} & \TBD{正式S/R}\\
    充电量（kWh） & \TBD{正式S/R} & \TBD{正式S/R} & \TBD{正式S/R} & \TBD{正式S/R} & \TBD{正式S/R}\\
    充电时长（min） & \TBD{正式S/R} & \TBD{正式S/R} & \TBD{正式S/R} & \TBD{正式S/R} & \TBD{正式S/R}\\
    会话最高SOC（\%） & \TBD{正式S/R} & \TBD{正式S/R} & \TBD{正式S/R} & \TBD{正式S/R} & \TBD{正式S/R}\\
    充电电费（元） & \TBD{正式S/R} & \TBD{正式S/R} & \TBD{正式S/R} & \TBD{正式S/R} & \TBD{正式S/R}\\
    单目标总账（元） & \TBD{正式S/R} & \TBD{正式S/R} & \TBD{正式S/R} & \TBD{正式S/R} & \TBD{正式S/R}\\
    充电排放（kgCO$_2$e） & \TBD{正式S/R} & \TBD{正式S/R} & \TBD{正式S/R} & \TBD{正式S/R} & \TBD{正式S/R}\\
    \bottomrule
  \end{tabular}
  }
  \tabnote{注：S表示简化充电函数下的账面求解，R表示同一方案按参考非线性函数回算。若算法在收益递减段前主动停止，本表应如实呈现其最高SOC与时长，而不是把“不触及拐点”写成机制失效。}
\end{table}

\section{结论与管理启示}

本文建立了时变电网碳强度下的动态多车场混合车队多趟车辆路径模型，并用同一全口径排放公式连接充电时刻、车型指派和动态发车决策。算法方面给出了成熟混合遗传搜索底座上的串行改进框架，但公开算例稳定排位与被接受的增量式消融仍为\TBD{正式算法证据}。

目前可确认的充电时刻作用空间是：固定排班下，碳感知择时相对有空即充移动174.732~kWh，充电侧排放减少59.082~kgCO$_2$e；在0.075元/kg碳价下，完整目标增加22.094元，配对翻转点为0.449元/kg。车型层目前只有两档探针，两档实际派遣电动车均为0；动态层没有可用结果。因此三层的统一结论仍为\TBD{车型完整扫描与动态正式结果}。

协同效率批显示，两家单干成本之和3922.079元，联合成本3033.206元，节约888.873元；同一联合解事后分摊后，两家分别节约444.436元，参与条件由1/2变为2/2满足。这一结果分别回答“联合运营是否节约”和“节约能否让各方都受益”，不把公平约束重新塞入效率搜索。

管理上，现有数据表明提高碳核算时间分辨率能够揭示充电时刻的排放差异，但现行价格信号不足以覆盖电费代价。对运营方，应同时观察车辆的真实可充窗口和电油车型的实际派遣；对政策制定者，需联合审视分时电价与碳价信号。关于公开算法性能、车型结构转移、动态需求响应及非线性充电的最终管理含义，均保留为\TBD{对应被接受实验}。

\begin{thebibliography}{99}
\bibitem{ref:state-council-2021}
国务院. 2030年前碳达峰行动方案[EB/OL]. 2021-10-26.

\bibitem{ref:li-2024}
Li Z, Chen Z, Li H, et al. On the value of orderly electric vehicle charging in carbon emission reduction[J]. Transportation Research Part D: Transport and Environment, 2024, 135: 104383.

\bibitem{ref:cheng-2022}
Cheng K W, Bian Y, Shi Y, et al. Carbon-aware EV charging[C]. 2022 IEEE SmartGridComm, 2022: 186--192.

\bibitem{ref:goeke-2015}
Goeke D, Schneider M. Routing a mixed fleet of electric and conventional vehicles[J]. European Journal of Operational Research, 2015, 245(1): 81--99.

\bibitem{ref:li-decheng-2021}
李得成, 陈彦如, 张宗成. 基于分支定价算法的电动车与燃油车混合车辆路径问题研究[J]. 系统工程理论与实践, 2021, 41(4): 995--1009.

\bibitem{ref:chen-wanru-2023}
陈婉茹, 徐光明, 张得志, 等. 碳交易机制下多中心混合车队配送路径和速度优化研究[J]. 系统工程理论与实践, 2023, 43(11): 3320--3335.

\bibitem{ref:gendreau-1999}
Gendreau M, Guertin F, Potvin J Y, et al. Parallel tabu search for real-time vehicle routing and dispatching[J]. Transportation Science, 1999, 33(4): 381--390.

\bibitem{ref:jiang-2024}
姜广田, 纪皎月, 董佳伟. 绿色物流配送下的多车型动态车辆路径优化[J]. 系统工程理论与实践, 2024, 44(7): 2362--2380.

\bibitem{ref:zhen-2020}
Zhen L, Ma C, Wang K, et al. Multi-depot multi-trip vehicle routing problem with time windows and release dates[J]. Transportation Research Part E, 2020, 135: 101866.

\bibitem{ref:froger-2022}
Froger A, Jabali O, Mendoza J E, et al. The electric vehicle routing problem with capacitated charging stations[J]. Transportation Science, 2022, 56(2): 460--482.

\bibitem{ref:keskin-2016}
Keskin M, Çatay B. Partial recharge strategies for the electric vehicle routing problem with time windows[J]. Transportation Research Part C, 2016, 65: 111--127.

\bibitem{ref:vidal-2012}
Vidal T, Crainic T G, Gendreau M, et al. A hybrid genetic algorithm for multidepot and periodic vehicle routing problems[J]. Operations Research, 2012, 60(3): 611--624.

\bibitem{ref:wouda-2024}
Wouda N A, Lan L, Kool W. PyVRP: A high-performance VRP solver package[J]. INFORMS Journal on Computing, 2024, 36(4): 943--955.

\end{thebibliography}

\end{document}
